\section{Evaluation}
\label{sec:eval}

We ask four questions.
\begin{itemize}
  \item[\rqref{1}] Does \tool find previously unknown defects in an engine that
        already fuzzes itself, and are they accepted upstream?
  \item[\rqref{2}] At an equal budget, which defects does each source of
        statements reach: mined templates, the typed generator, or both?
  \item[\rqref{3}] Do the defects lie in the blind spot of the engine's own
        fuzzers, or could those fuzzers have found them with more time?
  \item[\rqref{4}] What does \tool cost in time and in requests?
\end{itemize}

\subsection{Subjects}
\label{sec:subjects}
We test one engine and mine one other. We chose the engine under test with the
following criteria:
\begin{enumerate}
  \item it is actively maintained and reviews outside contributions, so a report
        gets an answer within weeks;
  \item it can be built from source at any commit, so a defect can be confirmed
        on the current main branch and shown fixed afterwards;
  \item it already runs its own fuzzers, so a defect we find is one those fuzzers
        missed rather than one nobody looked for.
\end{enumerate}
Turso meets all three. We used version 0.9.0-pre.2 at four commits: 730017f96,
where the first campaign ran; b588fbdc6, the last commit before our fixes;
ec86571642a6, which has only our first patch; and 7f0a9e6f4, which has both.
We mined DuckDB because it labels confirmed issues and because its fixes usually
carry a regression test with the triggering query~\cite{duckdb27846}. From 600
closed issues \tool kept 200, and from those it extracted 659 reproductions: 526
from regression tests and 133 from issue bodies.

\subsection{Baselines}
\label{sec:baselines}
\begin{description}
  \item[Mined only.] \tool with templates from DuckDB reports and no generator.
  \item[Generated only.] \tool with the typed generator and no templates. This is
        the plain baseline: a grammar-based fuzzer with boundary values.
  \item[Both.] Each case is drawn from either source.
  \item[Turso's own fuzzers.] Not run by us. We read their source to answer
        \rqref{3}.
\end{description}
All three runs use the same budget of 20{,}000 cases, the same seeds 0 to
19{,}999, the same crash oracle, and the same binary, b588fbdc6, which has none
of our patches.

\subsection{Measurements}
\begin{description}
  \item[\rqref{1}.] Defects found, their upstream state, the exit status and
        output of each reduced statement before and after our patches, and
        whether the tool regenerates each one from a single seed.
  \item[\rqref{2}.] Per run: cases executed, the share the engine accepted, oracle
        hits, and distinct signatures, with the first seed that reaches each.
        The same runs repeated on the build with only our first patch.
  \item[\rqref{3}.] For each defect, whether its function is in Turso's fuzzer
        vocabulary, whether its literals are in that fuzzer's domain, and whether
        SQLite implements the function.
  \item[\rqref{4}.] Wall-clock time and cases per second for each run, and
        requests made to build the corpus.
\end{description}

\subsection{Setup}
All measurements ran on one laptop: an Intel Core i7-1360P with 16 logical CPUs
and 15\,GB of memory, under Linux kernel 7.0, with Python 3.14.4 and Rust
1.88.0. Each case had a 10-second time limit. The statement for defect~4 also ran
under a 6\,GB memory limit, so the pre-fix build would fail fast instead of
building a string of several gigabytes. Each run in Section~\ref{sec:rq2} was
done once; Section~\ref{sec:threats} discusses what that allows us to claim.

\subsection{\rqref{1}: Does \tool Find New Defects?}
\label{sec:rq1}

\begin{table}[t]
  \caption{The four defects, each run as one statement in a fresh process.
  \emph{Before}: Turso b588fbdc6, which has none of our patches.
  \emph{After}: 7f0a9e6f4, which has both. Exit status 134 is SIGABRT.
  \emph{Found by}: how the defect was first found. \emph{Seed}: the seed that
  makes \tool regenerate the defect in a single case in the run that draws from
  both sources. Defect~4 ran under a 6\,GB memory limit;
  3{,}999{,}999{,}996 is $10^9$ characters times 4 bytes, minus the
  one-character input.}
  \label{tab:defects}
  \footnotesize
  \begin{tabular}{@{}rllrrll@{}}
    \toprule
     & Function & Defect & Before & After & Found by & Seed \\
    \midrule
    1 & \sql{lpad}   & allocates before the length check & 134, 2.28\,s & 1, 0.02\,s & \tool & 4{,}553 \\
    2 & \sql{rpad}   & allocates before the length check & 134, 2.24\,s & 1, 0.02\,s & \tool & 101 \\
    3 & \sql{repeat} & allocates first; size can overflow & 134, 2.27\,s & 1, 0.02\,s & \tool & 5{,}080 \\
    4 & \sql{lpad}   & counts characters, limit is bytes & 134, 3.96\,s & 1, 0.02\,s & review of our patch & -- \\
    \bottomrule
  \end{tabular}
\end{table}

Table~\ref{tab:defects} lists the defects. Defects 1 to 3 share a cause:
\sql{lpad}, \sql{rpad} and \sql{repeat} sized their result from the caller's
argument and allocated it before comparing it with
\texttt{Value::MAX\_BLOB\_LENGTH}. At a length near \imax the allocation fails and
Rust aborts the process before the check, which ran only on the finished
result, is reached. \sql{lpad} and \sql{rpad} share one function body,
\texttt{pad()}, so defects 1 and 2 have one root cause. We count them separately
because each aborts the engine on its own, and a user who calls only \sql{rpad}
is exposed. \sql{repeat} is a separate function that also multiplies, so its fix
needed an overflow check as well. \tool found all three: the first abort came
eight seconds into the first campaign. One patch fixed
them~\cite{turso11219,turso11225}.

Defect 4 became reachable only after that patch, and \tool did not find it. With
the check now running first, it compared a count of characters with a limit in
bytes. A four-byte fill such as an emoji passes a check of $10^9$ characters, and
about 4\,GB is built before the byte limit rejects it. We noticed the unit
mismatch while writing the first patch, but our pull request described it as
``conservative'', meaning harmless. Turso's automated reviewer replied that it
was permissive and supplied a demonstrating input. We then diagnosed the root
cause, and a second patch computes the exact byte length before building
anything~\cite{turso11260,turso11259}. On a debug build without a memory limit,
the statement ran for about 34 seconds and peaked at 11.7\,GB.

\begin{observation}
\label{obs:found}
\tool found three previously unknown defects in Turso, each of which ends the
process from one statement. Fixing them exposed a fourth, which \tool did not
find. Our two patches, which fix all four, were merged upstream.
\end{observation}

\begin{observation}
\label{obs:seed}
\tool regenerates each of the three defects it found from one recorded seed, in
a single case, on the pre-fix build: seed 4{,}553 for defect 1, 101 for defect 2
and 5{,}080 for defect 3. We checked that each seed calls the function its
defect names.
\end{observation}

\subsection{\rqref{2}: Mined Templates versus Typed Generation}
\label{sec:rq2}

\begin{table}[t]
  \caption{Each source of statements at an equal budget of 20{,}000 cases and
  seeds 0 to 19{,}999. \emph{Valid}: share of executed cases the engine
  accepted. \emph{Hits}: oracle violations on the pre-fix build. \emph{Distinct}:
  distinct signatures, the unit of comparison, since eight hits on one signature are one defect found
  eight times. The next three columns give the first seed that reached each
  signature, or a dash if none did. Defects 1 and 2 share a signature; in all
  three runs its first hit was an \sql{rpad} call. The third signature,
  \emph{Defect 1 again}, is defect~1 reached with a length above $10^9$ computed from table data.
  \emph{After patch 1}: hits when the same run is repeated on ec86571642a6, the
  build with only our first patch.}
  \label{tab:ablation}
  \footnotesize
  \begin{tabular}{@{}lrrrrrrrr@{}}
    \toprule
    Source & Executed & Valid & Hits & Distinct & Defects 1, 2 & Defect 3 & Defect 1 again & After patch 1 \\
    \midrule
    Mined only     & 19{,}169 &  8.9\% &  8 & 2 & 914 & 10{,}180 & --       & 0 \\
    Generated only & 20{,}000 & 98.3\% & 12 & 3 & 762 & 5{,}367  & 6{,}376  & 0 \\
    Both           & 19{,}759 & 71.1\% & 15 & 3 & 101 & 5{,}080  & 5{,}429  & 0 \\
    \bottomrule
  \end{tabular}
\end{table}

Table~\ref{tab:ablation} gives the result. Both sources reached both root causes
of defects 1 to 3: the shared \texttt{pad()} body and \sql{repeat}. The generator
reached one more signature. On inspection it is defect~1 again: an \sql{lpad}
call whose length, above $10^9$, is computed from table data. We first took this
signature for defect~4. Repeating every run on the build with only our first
patch settled it: all hits disappear, while that build still aborts on the
defect-4 statement. So no source reached defect~4.

The reason is in the boundary pools. Defect~4 needs a length of at most $10^9$
characters whose result is more than $10^9$ bytes, so between $2.5\times10^8$ and
$10^9$ for a four-byte fill. The integer pools jump from 32{,}767 to
2{,}147{,}483{,}647. None of the 60{,}000 cases that the three modes generate for
seeds 0 to 19{,}999 contains an integer literal in that range. The pools hold the limits of machine integers, not the
limits the engine sets itself.

Mined templates are wasteful but well aimed. Only 8.9\% of their cases ran at all,
because most DuckDB reproductions use syntax Turso rejects. Still, those few
valid cases reached both root causes. Drawing from both sources found the first
defect at seed 101, against 762 and 914 for either source alone, and recorded the
most hits.

\begin{observation}
\label{obs:gen}
At an equal budget, mined templates and typed generation each reached both root
causes of defects 1 to 3. No source reached defect~4, because no boundary pool
holds a length near the engine's own limit of $10^9$.
\end{observation}

\begin{observation}
\label{obs:both}
Drawing from both sources reached the first defect at seed 101, compared with 762
and 914 for each source alone, and recorded 15 hits against 12 and 8.
\end{observation}

\subsection{\rqref{3}: Could Turso's Own Fuzzers Have Found Them?}
\label{sec:rq3}

\begin{table}[t]
  \caption{Why Turso's own fuzzers cannot find the four defects, measured on the
  source at b588fbdc6~\cite{tursofuzzer}. \emph{In vocabulary}: the function is among the
  47 scalar functions of the differential fuzzer, or among those its four
  libFuzzer targets call. \emph{Literal in domain}: the defect's literal lies in the fuzzer's
  integer range $[-10^6, 10^6]$ or its alphanumeric string set. \emph{Judgeable}:
  SQLite implements the function.}
  \label{tab:blind}
  \footnotesize
  \begin{tabular}{@{}rllccc@{}}
    \toprule
     & Function & Literal the defect needs & In vocabulary & Literal in domain & Judgeable \\
    \midrule
    1 & \sql{lpad}   & integer $9.2\times10^{18}$ & \xmark & \xmark & \xmark \\
    2 & \sql{rpad}   & integer $9.2\times10^{18}$ & \xmark & \xmark & \xmark \\
    3 & \sql{repeat} & integer $9.2\times10^{18}$ & \xmark & \xmark & \xmark \\
    4 & \sql{lpad}   & a multi-byte fill character  & \xmark & \xmark & \xmark \\
    \bottomrule
  \end{tabular}
\end{table}

Table~\ref{tab:blind} answers the question with three independent checks. The
differential fuzzer's table of 47 scalar functions contains none of
\sql{lpad}, \sql{rpad} or \sql{repeat}. None of its four libFuzzer targets, and
no Rust file under its \texttt{testing} directory, mentions them either. The
functions are covered only by hand-written unit tests with small values. Even if
the functions were added, the default integer range stops at $10^6$, and defects
1 to 3 need a length near \imax. Strings are alphanumeric, and defect~4 needs a
multi-byte character. Finally, SQLite has none of the three functions, so a
differential oracle would have nothing to compare against.

\begin{observation}
\label{obs:blind}
All four defects lie in the blind spot of Turso's own fuzzers on every axis we
checked: the functions are absent from their vocabulary, the required literals
are outside their domain, and the reference engine lacks the functions. More time
would not have found them.
\end{observation}

\subsection{\rqref{4}: What Does \tool Cost?}
\label{sec:rq4}

\begin{table}[t]
  \caption{Cost of each 20{,}000-case run in Table~\ref{tab:ablation}, on the
  pre-fix build and the laptop described in the setup.
  \emph{Time}: wall-clock seconds for the whole run. \emph{Rate}: cases executed
  per second. \emph{First defect}: the case number of the first oracle hit; the run that draws
  from both sources is the slowest per case and the fastest to its first defect.}
  \label{tab:cost}
  \footnotesize
  \begin{tabular}{@{}lrrr@{}}
    \toprule
    Source & Time (s) & Rate (cases/s) & First defect (case) \\
    \midrule
    Mined only     & 378.0 & 50.7 & 914 \\
    Generated only & 388.8 & 51.4 & 762 \\
    Both           & 448.3 & 44.1 & 101 \\
    \bottomrule
  \end{tabular}
\end{table}

Table~\ref{tab:cost} gives the cost of each run. Every run finished in under eight
minutes. Drawing from both sources was the slowest per case and the fastest to a
defect: at 44.1 cases per second, case 101 is reached in about two seconds. We do
not have an explanation for the lower rate, and with one run per source we cannot
tell it apart from noise. Building the corpus took 387 requests to the GitHub API,
354 of them answered from a local cache, and is done once. Step~1 costs reading
time rather than compute. The reducer shrank the case for defect~3 to a single
statement in five executions.

\begin{observation}
\label{obs:cost}
One run of 20{,}000 cases took under eight minutes on a laptop. With both sources,
the first defect appeared at case 101, about two seconds into the run.
\end{observation}
